\documentclass[12pt, a4paper]{article}

% --- PACOTES ---
\usepackage[utf8]{inputenc}
\usepackage[brazil]{babel}
\usepackage{abntex2cite}
\usepackage{graphicx}
\usepackage{geometry}
\usepackage{hyperref}
\usepackage{indentfirst}
\usepackage{setspace}
\usepackage{times}
\usepackage{tikz}
\usepackage{array}
\usepackage{listings}
\usepackage{xcolor}
\usepackage{float}  % Necessário para a opção [H] em figures
\usetikzlibrary{shapes.geometric, arrows, positioning, fit, backgrounds}

% Configuração de listagens de código
\lstset{
    basicstyle=\ttfamily\footnotesize,
    breaklines=true,
    frame=single,
    numbers=left,
    numberstyle=\tiny,
    backgroundcolor=\color{gray!10},
    extendedchars=true,
    inputencoding=utf8,
    literate=%
    {á}{{\'a}}1 {é}{{\'e}}1 {í}{{\'i}}1 {ó}{{\'o}}1 {ú}{{\'u}}1
    {Á}{{\'A}}1 {É}{{\'E}}1 {Í}{{\'I}}1 {Ó}{{\'O}}1 {Ú}{{\'U}}1
    {à}{{\`a}}1 {è}{{\`e}}1 {ì}{{\`i}}1 {ò}{{\`o}}1 {ù}{{\`u}}1
    {À}{{\`A}}1 {È}{{\'E}}1 {Ì}{{\`I}}1 {Ò}{{\`O}}1 {Ù}{{\`U}}1
    {ã}{{\~a}}1 {õ}{{\~o}}1 {Ã}{{\~A}}1 {Õ}{{\~O}}1
    {â}{{\^a}}1 {ê}{{\^e}}1 {î}{{\^i}}1 {ô}{{\^o}}1 {û}{{\^u}}1
    {Â}{{\^A}}1 {Ê}{{\^E}}1 {Î}{{\^I}}1 {Ô}{{\^O}}1 {Û}{{\^U}}1
    {ç}{{\c{c}}}1 {Ç}{{\c{C}}}1
}

% --- MARGENS / FORMATAÇÃO ---
\geometry{
 a4paper,
 left=3cm,
 right=2cm,
 top=3cm,
 bottom=2cm
}
\onehalfspacing
\setlength{\parindent}{1.25cm}

\hypersetup{
    colorlinks=true,
    linkcolor=black,
    urlcolor=blue,
    citecolor=black
}

\begin{document}

% --- CAPA ---
\begin{titlepage}
    \centering
    {\large PROJETO FÍSICO \par}
    \vspace{1.5cm}
    {\bfseries\Large ARQUITETURA EFICIENTE PARA ESTIMAÇÃO DE POSE COMPLETA (\textit{FULL-BODY}) EM AMBIENTE VEICULAR ATRAVÉS DE \textit{LIFTING} 2D-PARA-3D \par}
    \vspace{2.5cm}

    \begin{flushleft}
        \textbf{Integrante:} Davi Baechtold Campos \\[0.3cm]
        \textbf{Orientador:} Prof. Dr. Alceu de Souza Brito Junior \\
        \textbf{Coorientador:} Prof. Dr. Alessandro Zimmer
    \end{flushleft}

    \vfill
    \begin{flushleft}
        \textbf{Visto do Orientador:}\\[1.4cm]
        \makebox[9cm]{\hrulefill}\\
        \small Prof. Dr. Alceu de Souza Brito Junior \\[0.5cm]
        Data: \underline{\hspace{3cm}}
    \end{flushleft}

    \vfill
    {\large Curitiba, \today \par}
\end{titlepage}

\tableofcontents
\newpage
\setcounter{page}{1}

% --- RESUMO ---
\section{RESUMO}

A segurança veicular representa um desafio crítico na indústria automotiva moderna, com sistemas avançados de proteção aos ocupantes (airbags, pré-tensionadores) demandando informações precisas sobre a pose tridimensional dos passageiros para otimizar sua atuação. Este projeto propõe o desenvolvimento de uma arquitetura eficiente para estimação de pose completa (\textit{full-body}) em ambiente veicular utilizando técnicas de \textit{2D-to-3D lifting} baseadas em \textit{keypoints}. O termo \textit{lifting} refere-se ao processo de inferir coordenadas tridimensionais (3D) a partir de pontos bidimensionais (2D) detectados em imagens, resolvendo a ambiguidade de profundidade inerente à projeção monocular através de modelos de aprendizado profundo que incorporam \textit{priors} anatômicos e geométricos.

A solução proposta viabiliza o uso de câmeras monoculares infravermelhas no desenvolvimento de métodos robustos para detecção e \textit{lifting} de \textit{keypoints} corporais, com expansão posterior incremental para face e mãos. A pesquisa aborda a questão fundamental de como construir uma arquitetura computacionalmente eficiente que integre corpo, face e mãos em um sistema unificado de monitoramento veicular. A arquitetura modular proposta compreende quatro subsistemas especializados: (1) módulo de aquisição e pré-processamento de imagens infravermelhas com correção de distorção e normalização adaptativa; (2) rede neural de detecção 2D multi-head baseada em RTMPose, adaptada para imagens grayscale e capaz de detectar simultaneamente 17 keypoints corporais, 68 faciais e 42 das mãos (21 por mão); (3) módulo de fusão temporal e lifting 3D empregando arquitetura transformer com atenção espaço-temporal sobre janelas de 16 frames, capaz de recuperar a profundidade e suavizar ruídos temporais; e (4) sistema de visualização em tempo real e cálculo de métricas de validação (MPJPE, PA-MPJPE, PCK).

O desenvolvimento segue metodologia incremental em seis fases distribuídas ao longo de 11 meses: Fase 1 (atual) estabelece o detector 2D corporal e pipeline de treinamento; Fase 2 implementa o lifting 3D com fusão temporal; Fase 3 consolida os resultados através de publicação científica; Fases 4-5 expandem para face e mãos; Fase 6 realiza otimizações finais e documentação. O projeto utilizará pré-treinamento com datasets públicos amplamente reconhecidos, incluindo o dataset COCO para keypoints 2D iniciais, Human3.6M para lifting 3D corporal, e o dataset do RTMPose como referência para arquiteturas de detecção 2D. A validação final será realizada no Drive\&Act dataset, específico para contexto automotivo, contendo imagens infravermelhas de motoristas em condições reais (variação de iluminação diurna/noturna, oclusões por componentes veiculares, movimentos naturais de condução). As métricas-alvo estabelecidas são: MPJPE $<$ 50mm para pose corporal, PCK@50mm $>$ 85\%, e taxa de processamento $\geq$ 20 FPS em GPU NVIDIA RTX 3060, garantindo viabilidade para aplicações em tempo real. A contribuição científica deste trabalho reside na adaptação e integração de técnicas estado-da-arte para o domínio específico veicular com imagens infravermelhas, área ainda pouco explorada na literatura, fornecendo uma solução viável economicamente para democratização de sistemas avançados de segurança.

\vspace{0.5cm}
\noindent\textbf{Palavras-chave:} Estimação de Pose 3D; Deep Learning; Visão Computacional; Sistemas Veiculares; Câmeras Infravermelhas; Monitoramento de Ocupantes; Segurança Automotiva.

% --- INTRODUÇÃO ---
\section{INTRODUÇÃO}

\subsection{Contexto e Motivação}

A indústria automotiva global enfrenta um paradoxo crítico: enquanto os avanços tecnológicos em sistemas de segurança veicular reduziram significativamente as fatalidades em acidentes nas últimas décadas, os sistemas de proteção aos ocupantes ainda operam de forma majoritariamente reativa e não-adaptativa. Segundo dados da Organização Mundial da Saúde (OMS), aproximadamente 1,35 milhão de pessoas morrem anualmente em acidentes de trânsito, e entre 20 a 50 milhões sofrem lesões não fatais. Sistemas de segurança passiva, como airbags e pré-tensionadores de cinto, são projetados para cenários padronizados de ocupação, assumindo posições corporais normativas que frequentemente não correspondem à realidade.

A próxima geração de sistemas de segurança veicular demanda capacidade de percepção contextual: conhecer em tempo real a pose tridimensional completa dos ocupantes (posição do tronco, orientação da cabeça, posição das mãos) para adaptar dinamicamente parâmetros críticos como força de acionamento de airbags, tensão dos cintos de segurança e até mesmo decisões de frenagem emergencial em veículos autônomos. Um ocupante inclinado para frente, por exemplo, requer acionamento diferenciado do airbag comparado a um ocupante reclinado; uma criança mal posicionada pode sofrer lesões graves por um airbag convencional; um motorista distraído manipulando objetos necessita alertas específicos.


\subsection{Problema de Pesquisa}

O problema central abordado neste trabalho pode ser formulado da seguinte maneira: \textit{Como estimar com precisão submilimétrica (MPJPE $<$ 50mm) a pose tridimensional completa (corpo, face e mãos) de ocupantes veiculares em tempo real ($\geq$ 20 FPS) utilizando exclusivamente câmeras monoculares infravermelhas de baixo custo, em condições desafiadoras de iluminação variável, oclusões parciais e posturas não-convencionais?}

Este problema se desdobra em múltiplos desafios técnicos inter-relacionados:

\textbf{Desafio 1 --- Adaptação de Domínio:} Modelos estado-da-arte de estimação de pose são treinados majoritariamente em datasets RGB de ambientes abertos (COCO, MPII, Human3.6M). A transferência para imagens infravermelhas implica em \textit{domain shift} significativo: ausência de informação cromática, texturas visuais distintas, padrões de ruído diferenciados. Estudos prévios demonstram degradação de 15-30\% em acurácia quando modelos RGB são aplicados diretamente a imagens grayscale sem retreinamento.

\textbf{Desafio 2 --- Especificidades do Ambiente Veicular:} O interior de veículos apresenta restrições espaciais severas, iluminação altamente variável (dia, noite, túneis, sombras), e oclusões sistemáticas causadas por componentes estruturais (volante obstruindo tronco e mãos, painel limitando visibilidade de pernas, cintos de segurança cruzando o corpo). Posturas típicas de condução (braços estendidos ao volante, rotação de cabeça para verificação de espelhos) diferem substancialmente das poses presentes em datasets genéricos.

\textbf{Desafio 3 --- Integração Full-Body:} A maioria dos trabalhos em estimação de pose foca isoladamente em corpo (17-21 keypoints), face (68-98 keypoints) ou mãos (21 keypoints por mão). Sistemas que integram todas as modalidades simultaneamente são escassos, especialmente em configurações de tempo real. A integração requer balanceamento cuidadoso de resolução espacial (face demanda alta resolução, corpo opera em baixa resolução), capacidade computacional e precisão por região anatômica.

\textbf{Desafio 4 --- Ambiguidade Monocular 2D-3D:} Câmeras monoculares projetam o espaço tridimensional em um plano bidimensional, resultando em perda irreversível de informação de profundidade. Múltiplas configurações 3D distintas podem produzir projeções 2D idênticas. Técnicas de \textit{lifting} devem aprender \textit{priors} anatômicos e geométricos (comprimentos de membros, ângulos articulares fisiologicamente plausíveis) para resolver essas ambiguidades de forma robusta.

\textbf{Desafio 5 --- Restrições de Tempo Real:} Aplicações de segurança veicular demandam latência inferior a 100ms para viabilizar intervenções preventivas. Considerando processamento a 20 FPS, o orçamento computacional é de apenas 50ms por frame, incluindo aquisição, detecção 2D, lifting 3D e pós-processamento.

\subsection{Objetivos}

\subsubsection{Objetivo Geral}

Desenvolver e validar uma arquitetura computacional eficiente para estimação de pose tridimensional completa (\textit{full-body}) de ocupantes veiculares em tempo real, utilizando câmeras monoculares infravermelhas de baixo custo, através de técnicas de detecção 2D de keypoints e lifting 2D-para-3D baseado em redes neurais profundas.

\subsubsection{Objetivos Específicos}

\begin{enumerate}
    \item \textbf{Adaptar e treinar} um modelo de detecção 2D de keypoints corporais baseado em RTMPose para operar robustamente em imagens infravermelhas (\textit{grayscale}), alcançando PCK@50mm superior a 85\% no dataset Drive\&Act;
    
    \item \textbf{Implementar} uma rede neural de lifting 2D-para-3D baseada em arquitetura transformer com fusão temporal sobre janelas deslizantes de 16 frames, capaz de estimar coordenadas 3D com MPJPE inferior a 50mm para pose corporal;
    
    \item \textbf{Expandir} a arquitetura para detecção integrada de face (68 keypoints) e mãos (21 keypoints por mão) através de heads especializados em rede multi-task, mantendo desempenho de tempo real;
    
    \item \textbf{Validar} experimentalmente o sistema completo no dataset Drive\&Act, analisando degradação de desempenho em condições adversas (oclusões, variação de iluminação, movimentos rápidos) e comparando com baselines da literatura;
    
    \item \textbf{Otimizar} a arquitetura para inferência em tempo real ($\geq$ 20 FPS) em hardware de gama média (GPU NVIDIA RTX 3060), explorando técnicas de quantização, pruning e otimização de operadores;
    
    \item \textbf{Documentar} e disponibilizar código-fonte, modelos treinados e pipeline de treinamento/avaliação para reprodutibilidade científica e fomento de pesquisas futuras na área.
\end{enumerate}

\subsection{Contribuições Esperadas}

Este trabalho visa contribuir para o estado-da-arte em múltiplas dimensões:

\begin{itemize}
    \item \textbf{Contribuição Científica:} Primeira arquitetura integrada full-body (corpo + face + mãos) especificamente adaptada para imagens infravermelhas veiculares, com análise sistemática de transferência de domínio RGB-para-grayscale;
    
    \item \textbf{Contribuição Técnica:} Pipeline completo de treinamento, validação e inferência otimizado para tempo real, incluindo estratégias de data augmentation específicas para contexto automotivo;
    
    \item \textbf{Contribuição Aplicada:} Demonstração de viabilidade técnica e econômica de sistemas de monitoramento de ocupantes baseados em hardware de baixo custo, potencialmente democratizando tecnologias de segurança avançadas;
    
    \item \textbf{Contribuição Metodológica:} Análise detalhada de riscos, validação incremental por módulos, e documentação extensiva que pode servir como referência para projetos similares em visão computacional automotiva.
\end{itemize}

\subsection{Escopo e Delimitações}

\subsubsection{Escopo Incluído}

\textbf{O que está incluído no projeto:}
\begin{itemize}
    \item Desenvolvimento completo de pipeline de detecção 2D de keypoints corporais, faciais e das mãos em imagens infravermelhas (\textit{grayscale});
    \item Implementação de módulo de lifting 2D-para-3D com fusão temporal baseado em transformers;
    \item Integração de detecção full-body em arquitetura unificada multi-head;
    \item Treinamento e fine-tuning de modelos em datasets públicos (COCO, Human3.6M, 300W, FreiHAND, Drive\&Act);
    \item Validação experimental extensiva com análise de degradação em condições adversas;
    \item Otimização para inferência em tempo real em GPU de gama média;
    \item Implementação de sistema de visualização 3D interativa e cálculo automático de métricas;
    \item Documentação técnica completa e disponibilização de código open-source.
\end{itemize}

\subsubsection{Escopo Excluído}

\textbf{O que NÃO está incluído no projeto:}
\begin{itemize}
    \item Desenvolvimento de hardware customizado (câmeras, processadores embarcados ou sistemas de montagem veicular);
    \item Implementação de sistemas completos de decisão de acionamento de airbags ou integração com ECUs automotivas;
    \item Coleta extensiva de novos datasets com ocupantes reais (utilizaremos exclusivamente datasets públicos existentes por questões éticas e logísticas);
    \item Implementação de modelos para detecção e rastreamento de múltiplos ocupantes simultâneos (foco exclusivo no motorista);
    \item Deploy e validação em sistemas embarcados automotivos reais ou veículos em operação (validação restrita a ambiente desktop/laptop);
    \item Certificações ou homologações para uso em sistemas de segurança automotivos reais (ISO 26262, ASPICE);
    \item Análise de privacidade, aspectos legais ou regulatórios de sistemas de monitoramento de ocupantes.
\end{itemize}

\subsection{Organização do Documento}

Este documento está estruturado de forma a apresentar progressivamente todos os aspectos técnicos, metodológicos e gerenciais do projeto:

\begin{itemize}
    \item \textbf{Seção 3 --- Detalhamento do Projeto:} Apresenta a arquitetura completa do sistema através de diagramas em blocos, especificação técnica detalhada de cada módulo (linguagens, bibliotecas, frameworks), fluxogramas de processamento, pseudocódigo de componentes críticos, e definição precisa de interfaces entre módulos;
    
    \item \textbf{Seção 4 --- Cronograma:} Detalha o planejamento temporal em 6 fases distribuídas ao longo de 11 meses, com marcos de validação, entregas parciais e status atual de cada atividade;
    
    \item \textbf{Seção 5 --- Procedimentos de Teste e Validação:} Descreve metodologia de validação em perspectivas de caixa preta (testes funcionais end-to-end) e caixa branca (validação de componentes internos), incluindo métricas, datasets de teste e critérios de aceitação;
    
    \item \textbf{Seção 6 --- Análise de Riscos:} Identifica e quantifica riscos técnicos, de dados e de cronograma, com probabilidades estimadas, impactos potenciais, severidades calculadas e estratégias de mitigação específicas;
    
    \item \textbf{Seção 7 --- Conclusão:} Sintetiza a viabilidade do projeto, contribuições esperadas, estado atual de desenvolvimento e perspectivas futuras;
    
    \item \textbf{Seção 8 --- Referências Bibliográficas:} Lista completa de trabalhos científicos, frameworks, bibliotecas e datasets utilizados como fundamentação teórica e ferramental do projeto.
\end{itemize}
    \item Certificações ou homologações para uso em sistemas de segurança automotivos reais (ISO 26262, ASPICE);
    \item Análise de privacidade, aspectos legais ou regulatórios de sistemas de monitoramento de ocupantes.
\end{itemize}

\subsection{Organização do Documento}

% --- DETALHAMENTO DO PROJETO ---
\section{DETALHAMENTO DO PROJETO}

\subsection{Diagrama em Blocos - Visão Geral}

A Figura \ref{fig:blocos_geral} apresenta a arquitetura geral do sistema, dividida em quatro módulos principais que processam sequencialmente os dados desde a captura até a saída final.

\begin{figure}[H]
    \centering
    \begin{tikzpicture}[
            node distance=2.5cm,
            block/.style={rectangle, draw, thick, fill=blue!20, text width=3cm, text centered, minimum height=1.2cm, rounded corners},
            input/.style={rectangle, draw, thick, fill=green!20, text width=2.2cm, text centered, minimum height=0.9cm},
            output/.style={rectangle, draw, thick, fill=red!20, text width=2.2cm, text centered, minimum height=0.9cm},
            arrow/.style={thick,->,>=stealth}
        ]
        
        % Input
        \node[input] (cam) {Câmeras IR\\Monoculares};
        
        % Blocos
        \node[block, right of=cam, node distance=4cm] (mod1) {\textbf{Módulo 1:}\\Aquisição e\\Pré-processamento};
        \node[block, right of=mod1, node distance=4.5cm] (mod2) {\textbf{Módulo 2:}\\Detecção 2D\\Keypoints};
        \node[block, below of=mod2, node distance=2.5cm] (mod3) {\textbf{Módulo 3:}\\Fusão Temporal e\\Lifting 3D};
        \node[block, left of=mod3, node distance=4.5cm] (mod4) {\textbf{Módulo 4:}\\Visualização e\\Métricas};
        
        % Output
        \node[output, left of=mod4, node distance=4cm] (out) {Pose 3D\\Full-Body};
        
        % Arrows
        \draw[arrow] (cam) -- (mod1);
        \draw[arrow] (mod1) -- node[above] {\tiny Frames IR} (mod2);
        \draw[arrow] (mod2) -- node[right] {\tiny Keypoints 2D} (mod3);
        \draw[arrow] (mod3) -- node[below] {\tiny Pose 3D} (mod4);
        \draw[arrow] (mod4) -- (out);
        
        % Feedback loop
        \draw[arrow, dashed] (mod3.west) -- ++(-1,0) |- node[left, pos=0.25] {\tiny Temporal} (mod2.west);
        
    \end{tikzpicture}
    \caption{Diagrama em blocos da arquitetura geral do sistema.}
    \label{fig:blocos_geral}
\end{figure}

\subsection{Módulo 1: Aquisição e Pré-processamento}

\subsubsection{Descrição Funcional}

O Módulo 1 é responsável pela captura de frames das câmeras infravermelhas e preparação dos dados para processamento pelos módulos subsequentes. Este módulo opera como a interface entre o hardware de captura e o pipeline de processamento.

\subsubsection{Especificação Técnica}

\textbf{Linguagem:} Python 3.10+\\
\textbf{Bibliotecas principais:}
\begin{itemize}
    \item OpenCV 4.8+ (cv2) - Captura de vídeo e processamento de imagem
    \item NumPy 1.24+ - Operações numéricas e manipulação de arrays
    \item Threading/Multiprocessing - Gerenciamento de múltiplas câmeras
\end{itemize}

\textbf{Hardware alvo:} Sistema desktop/laptop com GPU NVIDIA (RTX 3060 ou superior)\\
\textbf{Entrada:} Stream de vídeo de câmeras USB infravermelhas (resolução 640x480 a 30fps)\\
\textbf{Saída:} Frames normalizados (tensores PyTorch) com shape $[B, 1, H, W]$

\subsubsection{Diagrama de Fluxo}

\begin{figure}[H]
    \centering
    \begin{tikzpicture}[
            node distance=1.5cm,
            process/.style={rectangle, draw, fill=blue!15, text width=4cm, text centered, minimum height=0.8cm},
            decision/.style={diamond, draw, fill=yellow!20, text width=2.5cm, text centered, aspect=2},
            arrow/.style={thick,->,>=stealth}
        ]
        
        \node[process] (init) {Inicializar captura\\de câmera(s)};
        \node[process, below of=init] (cap) {Capturar frame};
        \node[decision, below of=cap, node distance=2cm] (valid) {Frame\\válido?};
        \node[process, below of=valid, node distance=2cm] (undist) {Corrigir distorção};
        \node[process, below of=undist] (resize) {Redimensionar\\para 256x256};
        \node[process, below of=resize] (norm) {Normalizar\\$[0, 255] \to [0, 1]$};
        \node[process, below of=norm] (tensor) {Converter para\\tensor PyTorch};
        \node[process, below of=tensor] (buffer) {Adicionar ao buffer\\temporal};
        \node[process, below of=buffer] (output) {Enviar para\\Módulo 2};
        
        \draw[arrow] (init) -- (cap);
        \draw[arrow] (cap) -- (valid);
        \draw[arrow] (valid) -- node[right] {Sim} (undist);
        \draw[arrow] (valid.west) -- ++(-1.5,0) node[above, pos=0.5] {Não} |- (cap.west);
        \draw[arrow] (undist) -- (resize);
        \draw[arrow] (resize) -- (norm);
        \draw[arrow] (norm) -- (tensor);
        \draw[arrow] (tensor) -- (buffer);
        \draw[arrow] (buffer) -- (output);
        \draw[arrow] (output.east) -- ++(1.5,0) |- (cap.east);
        
    \end{tikzpicture}
    \caption{Fluxograma do Módulo 1 - Aquisição e Pré-processamento.}
    \label{fig:mod1_flow}
\end{figure}

\subsubsection{Pseudocódigo}

\begin{lstlisting}[language=Python, caption=Módulo 1 - Aquisição]
class CameraAcquisition:
    def __init__(self, camera_id=0, resolution=(640,480)):
        self.cap = cv2.VideoCapture(camera_id)
        self.cap.set(cv2.CAP_PROP_FRAME_WIDTH, resolution[0])
        self.cap.set(cv2.CAP_PROP_FRAME_HEIGHT, resolution[1])
        self.camera_matrix = load_calibration()
        self.temporal_buffer = deque(maxlen=16)  # 16 frames
        
    def preprocess_frame(self, frame):
        # 1. Correcao de distorcao
        frame = cv2.undistort(frame, self.camera_matrix, 
                              self.dist_coeffs)
        
        # 2. Redimensionar para entrada da rede
        frame = cv2.resize(frame, (256, 256), 
                          interpolation=cv2.INTER_LINEAR)
        
        # 3. Normalizar para [0, 1]
        frame = frame.astype(np.float32) / 255.0
        
        # 4. Converter para tensor PyTorch [1, H, W]
        tensor = torch.from_numpy(frame).unsqueeze(0)
        
        return tensor
    
    def get_batch(self, batch_size=1):
        frames = []
        for _ in range(batch_size):
            ret, frame = self.cap.read()
            if not ret:
                raise RuntimeError("Falha na captura")
            
            tensor = self.preprocess_frame(frame)
            frames.append(tensor)
            self.temporal_buffer.append(tensor)
        
        # Stack em batch [B, 1, H, W]
        batch = torch.stack(frames, dim=0)
        return batch, self.temporal_buffer
\end{lstlisting}

\subsubsection{Interface de Saída}

\textbf{Formato:} Dicionário Python com as seguintes chaves:
\begin{lstlisting}[language=Python]
{
    'frame': torch.Tensor,        # Shape: [B, 1, H, W]
    'temporal_buffer': deque,     # Ultimos 16 frames
    'timestamp': float,           # Timestamp de captura
    'metadata': {
        'camera_id': int,
        'resolution': tuple,
        'fps': float
    }
}
\end{lstlisting}

\subsection{Módulo 2: Detecção 2D de Keypoints}

\subsubsection{Descrição Funcional}

O Módulo 2 implementa a detecção de keypoints 2D utilizando uma arquitetura baseada em RTMPose adaptada para imagens grayscale. O módulo será desenvolvido incrementalmente: primeiro corpo (17 keypoints), depois face (68 keypoints), e finalmente mãos (21 keypoints cada).

\subsubsection{Especificação Técnica}

\textbf{Linguagem:} Python 3.10+\\
\textbf{Framework:} PyTorch 2.0+ com CUDA 11.8\\
\textbf{Arquitetura base:} RTMPose (adaptada)\\
\textbf{Bibliotecas principais:}
\begin{itemize}
    \item torch, torchvision - Framework de deep learning
    \item timm - Modelos pré-treinados (backbones)
    \item mmpose - Toolkit de pose estimation (referência)
\end{itemize}

\textbf{Ferramentas de desenvolvimento:}
\begin{itemize}
    \item PyTorch Lightning - Organização de treinamento
    \item Weights \& Biases (wandb) - Tracking de experimentos
    \item TensorBoard - Visualização de métricas
\end{itemize}

\textbf{Entrada:} Batch de frames $[B, 1, 256, 256]$\\
\textbf{Saída:} Keypoints 2D com shape $[B, N_{\mathrm{keypoints}}, 3]$ (x, y, confidence)

\subsubsection{Arquitetura da Rede Neural}

\begin{figure}[H]
    \centering
    \begin{tikzpicture}[
            node distance=1.9cm,
            layer/.style={rectangle, draw, thick, fill=blue!15, text width=3.5cm, text centered, minimum height=0.8cm},
            arrow/.style={thick,->,>=stealth}
        ]
        
        \node[layer] (input) {Input\\{[}B, 1, 256, 256{]}};
        \node[layer, below of=input] (stem) {Stem Conv\\{[}B, 64, 128, 128{]}};
        \node[layer, below of=stem] (backbone) {EfficientNet-B0\\Backbone\\{[}B, 1280, 8, 8{]}};
        \node[layer, below of=backbone] (neck) {FPN Neck\\{[}B, 256, 32, 32{]}};
        \node[layer, below of=neck] (head_body) {Body Head\\{[}B, 17, 64, 64{]}};
        \node[layer, right of=head_body, node distance=4.5cm] (head_face) {Face Head\\{[}B, 68, 128, 128{]}};
        \node[layer, below of=head_body] (head_hands) {Hands Head\\{[}B, 42, 64, 64{]}};
        \node[layer, below of=head_hands] (output) {Keypoints 2D\\{[}B, N, 3{]}};
        
        \draw[arrow] (input) -- (stem);
        \draw[arrow] (stem) -- (backbone);
        \draw[arrow] (backbone) -- (neck);
        \draw[arrow] (neck) -- (head_body);
        \draw[arrow] (neck.east) -- (head_face);
        \draw[arrow] (neck.south) -- (head_hands);
        \draw[arrow] (head_body) -- (output);
        \draw[arrow] (head_face.south) -- (output.east);
        \draw[arrow] (head_hands) -- (output);
        
    \end{tikzpicture}
    \caption{Arquitetura da rede neural de detecção 2D com heads especializados.}
    \label{fig:mod2_arch}
\end{figure}

\subsubsection{Pseudocódigo - Modelo}

\begin{lstlisting}[language=Python, caption=Módulo 2 - Detecção 2D]
import torch
import torch.nn as nn
import timm

class Keypoint2DDetector(nn.Module):
    def __init__(self, num_body_kpts=17, 
                 num_face_kpts=68, num_hand_kpts=21):
        super().__init__()
        
        # Backbone compartilhado (adaptado para grayscale)
        self.backbone = timm.create_model(
            'efficientnet_b0', 
            pretrained=True,
            in_chans=1,  # Grayscale input
            features_only=True
        )
        
        # Feature Pyramid Network
        self.fpn = FeaturePyramidNetwork(
            in_channels=[40, 112, 320],
            out_channels=256
        )
        
        # Heads especializados
        self.body_head = KeypointHead(
            in_channels=256, 
            num_keypoints=num_body_kpts,
            output_size=(64, 64)
        )
        
        self.face_head = KeypointHead(
            in_channels=256,
            num_keypoints=num_face_kpts,
            output_size=(128, 128)  # Maior resolucao
        )
        
        self.hands_head = KeypointHead(
            in_channels=256,
            num_keypoints=num_hand_kpts * 2,  # 2 maos
            output_size=(64, 64)
        )
    
    def forward(self, x, mode='body'):
        # Extracao de features multi-escala
        features = self.backbone(x)
        fpn_features = self.fpn(features)
        
        outputs = {}
        if mode in ['body', 'full']:
            outputs['body'] = self.body_head(fpn_features)
        if mode in ['face', 'full']:
            outputs['face'] = self.face_head(fpn_features)
        if mode in ['hands', 'full']:
            outputs['hands'] = self.hands_head(fpn_features)
        
        return outputs

class KeypointHead(nn.Module):
    def __init__(self, in_channels, num_keypoints, 
                 output_size):
        super().__init__()
        self.upsample = nn.Sequential(
            nn.ConvTranspose2d(in_channels, 128, 4, 2, 1),
            nn.BatchNorm2d(128),
            nn.ReLU(),
            nn.ConvTranspose2d(128, 64, 4, 2, 1),
            nn.BatchNorm2d(64),
            nn.ReLU(),
        )
        
        # Heatmap prediction
        self.heatmap_conv = nn.Conv2d(
            64, num_keypoints, 1
        )
        
        # Offset regression para precisao sub-pixel
        self.offset_conv = nn.Conv2d(
            64, num_keypoints * 2, 1
        )
    
    def forward(self, x):
        x = self.upsample(x)
        heatmaps = self.heatmap_conv(x)
        offsets = self.offset_conv(x)
        
        # Extrair coordenadas dos heatmaps
        keypoints = self.heatmap_to_keypoints(
            heatmaps, offsets
        )
        return keypoints
\end{lstlisting}

\subsubsection{Estratégia de Treinamento}

\textbf{Fase 1 - Corpo (Fase atual do TCC):}
\begin{enumerate}
    \item \textbf{Pré-treinamento}: COCO dataset convertido para grayscale (200k imagens)
    \item \textbf{Fine-tuning}: Drive\&Act dataset (subset de treinamento)
    \item \textbf{Hiperparâmetros}:
    \begin{itemize}
        \item Learning rate: 1e-3 com cosine annealing
        \item Batch size: 32
        \item Epochs: 100 (early stopping com patience=15)
        \item Optimizer: AdamW (weight decay=1e-4)
        \item Loss: Combined MSE (heatmaps) + L1 (offsets)
    \end{itemize}
    \item \textbf{Data Augmentation}:
    \begin{itemize}
        \item Random rotation (-30° a +30°)
        \item Random scale (0.8 a 1.2)
        \item Random brightness/contrast
        \item Simulação de oclusão (volante, painel)
    \end{itemize}
\end{enumerate}

\textbf{Fase 2 - Face (expansão futura):}
\begin{enumerate}
    \item Integrar dataset 300W convertido para grayscale
    \item Treinar face\_head mantendo backbone congelado
    \item Fine-tuning conjunto corpo + face
\end{enumerate}

\textbf{Fase 3 - Mãos (expansão futura):}
\begin{enumerate}
    \item Integrar FreiHAND dataset
    \item Treinar hands\_head
    \item Fine-tuning full-body unificado
\end{enumerate}

\subsection{Módulo 3: Fusão Temporal e Lifting 3D}

\subsubsection{Descrição Funcional}

O Módulo 3 recebe os keypoints 2D detectados e realiza duas operações principais: fusão temporal para consistência entre frames e lifting 2D-para-3D para estimar coordenadas tridimensionais. A arquitetura é baseada em transformers para capturar dependências espaciais e temporais.

\subsubsection{Especificação Técnica}

\textbf{Linguagem:} Python 3.10+\\
\textbf{Framework:} PyTorch 2.0+ com CUDA\\
\textbf{Arquitetura:} Transformer Encoder com atenção espaço-temporal\\
\textbf{Entrada:} Sequência de keypoints 2D $[T, N_{\mathrm{keypoints}}, 2]$ onde $T=16$ frames\\
\textbf{Saída:} Keypoints 3D $[N_{\mathrm{keypoints}}, 3]$ para o frame central

\subsubsection{Arquitetura do Modelo de Lifting}

\begin{figure}[H]
    \centering
    \begin{tikzpicture}[
            node distance=2.9cm,
            layer/.style={rectangle, draw, thick, fill=green!5, text width=4.5cm, text centered, minimum height=1.8cm},
            arrow/.style={thick,->,>=stealth}
        ]
        
        \node[layer] (input) {Keypoints 2D\\Sequência\\{[}T=16, N=17, 2{]}};
        \node[layer, below of=input] (embed) {Embedding Posicional\\+ Projeção\\{[}T, N, D=256{]}};
        \node[layer, below of=embed] (trans1) {Transformer\\Encoder Layer 1};
        \node[layer, below of=trans1] (trans2) {Transformer\\Encoder Layer 2};
        \node[layer, below of=trans2] (trans3) {Transformer\\Encoder Layer 3};
        \node[layer, below of=trans3] (pool) {Temporal Pooling\\(Frame Central)\\{[}N, D{]}};
        \node[layer, below of=pool] (mlp) {MLP Regressora\\{[}N, D{]} → {[}N, 3{]}};
        \node[layer, below of=mlp] (output) {Keypoints 3D\\{[}N, 3{]}};
        
        \draw[arrow] (input) -- (embed);
        \draw[arrow] (embed) -- (trans1);
        \draw[arrow] (trans1) -- (trans2);
        \draw[arrow] (trans2) -- (trans3);
        \draw[arrow] (trans3) -- (pool);
        \draw[arrow] (pool) -- (mlp);
        \draw[arrow] (mlp) -- (output);
        
    \end{tikzpicture}
    \caption{Arquitetura do modelo de lifting 3D com fusão temporal.}
    \label{fig:mod3_arch}
\end{figure}

\subsubsection{Pseudocódigo}

\begin{lstlisting}[language=Python, caption=Módulo 3 - Lifting 3D]
class Lifting3DNetwork(nn.Module):
    def __init__(self, num_keypoints=17, d_model=256, 
                 num_layers=3, num_heads=8, 
                 temporal_window=16):
        super().__init__()
        self.num_keypoints = num_keypoints
        self.temporal_window = temporal_window
        
        # Projecao de 2D para embedding space
        self.input_projection = nn.Linear(2, d_model)
        
        # Positional encoding (espacial + temporal)
        self.spatial_pe = nn.Parameter(
            torch.randn(1, num_keypoints, d_model)
        )
        self.temporal_pe = nn.Parameter(
            torch.randn(1, temporal_window, 1, d_model)
        )
        
        # Transformer encoder
        encoder_layer = nn.TransformerEncoderLayer(
            d_model=d_model,
            nhead=num_heads,
            dim_feedforward=d_model * 4,
            dropout=0.1,
            activation='gelu',
            batch_first=True
        )
        self.transformer = nn.TransformerEncoder(
            encoder_layer,
            num_layers=num_layers
        )
        
        # MLP para regressao 3D
        self.output_mlp = nn.Sequential(
            nn.Linear(d_model, 512),
            nn.ReLU(),
            nn.Dropout(0.1),
            nn.Linear(512, 256),
            nn.ReLU(),
            nn.Dropout(0.1),
            nn.Linear(256, 3)  # (x, y, z)
        )
        
    def forward(self, keypoints_2d_sequence):
        # keypoints_2d_sequence: [B, T, N, 2]
        B, T, N, _ = keypoints_2d_sequence.shape
        
        # Projetar para embedding space
        x = self.input_projection(keypoints_2d_sequence)
        # x: [B, T, N, D]
        
        # Adicionar positional encodings
        x = x + self.spatial_pe.unsqueeze(1)
        x = x + self.temporal_pe
        
        # Reshape para processar todos tokens juntos
        x = x.view(B, T * N, -1)
        
        # Transformer encoder
        x = self.transformer(x)  # [B, T*N, D]
        
        # Reshape back
        x = x.view(B, T, N, -1)
        
        # Temporal pooling - focar no frame central
        center_idx = T // 2
        x_center = x[:, center_idx, :, :]  # [B, N, D]
        
        # Regressao 3D
        keypoints_3d = self.output_mlp(x_center)  # [B, N, 3]
        
        return keypoints_3d

# Pipeline de Lifting
class LiftingPipeline:
    def __init__(self, model, device='cuda'):
        self.model = model.to(device)
        self.device = device
        self.temporal_buffer = deque(maxlen=16)
        
    def process(self, keypoints_2d):
        # keypoints_2d: [N, 2] para frame atual
        self.temporal_buffer.append(keypoints_2d)
        
        if len(self.temporal_buffer) < 16:
            # Buffer nao completo, usar padding
            padded = self._pad_sequence(self.temporal_buffer)
        else:
            padded = torch.stack(list(self.temporal_buffer))
        
        # Forward pass
        with torch.no_grad():
            keypoints_3d = self.model(
                padded.unsqueeze(0).to(self.device)
            )
        
        return keypoints_3d.squeeze(0).cpu()
\end{lstlisting}

\subsubsection{Estratégia de Treinamento}

\textbf{Dataset:} Human3.6M (pré-treinamento) + Drive\&Act (fine-tuning)\\
\textbf{Loss Function:} MPJPE (Mean Per Joint Position Error)
\begin{lstlisting}[language=Python]
def mpjpe_loss(pred, target, valid_mask):
    # pred, target: [B, N, 3]
    # valid_mask: [B, N] (1 se visivel, 0 se oculto)
    diff = pred - target
    distances = torch.norm(diff, dim=2)  # [B, N]
    
    # Aplicar mascara para ignorar keypoints ocultos
    masked_distances = distances * valid_mask
    loss = masked_distances.sum() / valid_mask.sum()
    return loss
\end{lstlisting}

\textbf{Hiperparâmetros:}
\begin{itemize}
    \item Learning rate: 5e-4 com ReduceLROnPlateau
    \item Batch size: 64 sequências
    \item Epochs: 80
    \item Optimizer: AdamW
    \item Weight decay: 1e-5
\end{itemize}

\subsection{Módulo 4: Visualização e Métricas}

\subsubsection{Descrição Funcional}

O Módulo 4 é responsável por apresentar os resultados de forma visual e calcular métricas de avaliação. Este módulo é crítico para desenvolvimento, debug e validação do sistema.

\subsubsection{Especificação Técnica}

\textbf{Linguagem:} Python 3.10+\\
\textbf{Bibliotecas:}
\begin{itemize}
    \item Matplotlib 3.7+ - Visualização 2D/3D
    \item Open3D 0.17+ - Renderização 3D interativa
    \item PyQt5 - Interface gráfica
    \item Plotly - Visualizações interativas web
\end{itemize}

\subsubsection{Funcionalidades Implementadas}

\begin{enumerate}
    \item \textbf{Overlay 2D}: Renderizar keypoints detectados sobre frame original
    \item \textbf{Visualização 3D}: Esqueleto 3D com rotação interativa
    \item \textbf{Métricas em Tempo Real}: FPS, latência, MPJPE
    \item \textbf{Heatmaps de Confiança}: Visualizar confiança de detecção
    \item \textbf{Comparação Ground Truth}: Overlay de predição vs anotação
\end{enumerate}

\subsubsection{Pseudocódigo - Interface Principal}

\begin{lstlisting}[language=Python, caption=Módulo 4 - Visualização]
class VisualizationModule:
    def __init__(self):
        self.fig_2d = plt.figure(figsize=(12, 6))
        self.ax_2d = self.fig_2d.add_subplot(121)
        self.ax_3d = self.fig_2d.add_subplot(122, 
                                             projection='3d')
        self.metrics = MetricsCalculator()
        
    def render_2d(self, frame, keypoints_2d, confidence):
        self.ax_2d.clear()
        self.ax_2d.imshow(frame, cmap='gray')
        
        # Desenhar keypoints
        for i, (x, y) in enumerate(keypoints_2d):
            color = self._confidence_to_color(confidence[i])
            self.ax_2d.scatter(x, y, c=color, s=50)
        
        # Desenhar conexoes (skeleton)
        for (start, end) in BODY_CONNECTIONS:
            x_vals = [keypoints_2d[start][0], 
                      keypoints_2d[end][0]]
            y_vals = [keypoints_2d[start][1], 
                      keypoints_2d[end][1]]
            self.ax_2d.plot(x_vals, y_vals, 'b-', 
                           linewidth=2)
    
    def render_3d(self, keypoints_3d):
        self.ax_3d.clear()
        
        # Desenhar skeleton 3D
        xs, ys, zs = keypoints_3d.T
        self.ax_3d.scatter(xs, ys, zs, c='r', s=50)
        
        for (start, end) in BODY_CONNECTIONS:
            self.ax_3d.plot(
                [xs[start], xs[end]],
                [ys[start], ys[end]],
                [zs[start], zs[end]],
                'b-', linewidth=2
            )
        
        self._set_3d_view_properties()
    
    def update_metrics(self, pred_3d, gt_3d, 
                       inference_time):
        mpjpe = self.metrics.compute_mpjpe(pred_3d, 
                                           gt_3d)
        fps = 1.0 / inference_time
        
        # Exibir metricas
        metrics_text = f"MPJPE: {mpjpe:.2f}mm\n"
        metrics_text += f"FPS: {fps:.1f}\n"
        metrics_text += f"Latency: {inference_time*1000:.1f}ms"
        
        self.ax_2d.text(0.02, 0.98, metrics_text,
                       transform=self.ax_2d.transAxes,
                       verticalalignment='top',
                       bbox=dict(boxstyle='round', 
                                facecolor='wheat', 
                                alpha=0.5))

class MetricsCalculator:
    def compute_mpjpe(self, pred, gt):
        # pred, gt: [N, 3]
        distances = np.linalg.norm(pred - gt, axis=1)
        return np.mean(distances)
    
    def compute_pa_mpjpe(self, pred, gt):
        # Procrustes-aligned MPJPE
        pred_aligned = self._procrustes_align(pred, gt)
        return self.compute_mpjpe(pred_aligned, gt)
    
    def compute_pck(self, pred, gt, threshold=50):
        # Percentage of Correct Keypoints
        distances = np.linalg.norm(pred - gt, axis=1)
        correct = (distances < threshold).sum()
        return 100.0 * correct / len(distances)
\end{lstlisting}

\subsubsection{Interface Gráfica}

\begin{figure}[H]
    \centering
    \begin{tikzpicture}[scale=0.8]
        % Main window
        \draw[thick] (0,0) rectangle (14,8);
        \node at (7, 8.5) {\textbf{Sistema de Estimação de Pose 3D}};
        
        % Left panel - Input video
        \draw[thick] (0.5,4.5) rectangle (6.5,7.5);
        \node at (3.5, 7.8) {\small Frame IR + Keypoints 2D};
        
        % Right panel - 3D visualization
        \draw[thick] (7.5,4.5) rectangle (13.5,7.5);
        \node at (10.5, 7.8) {\small Visualização 3D};
        
        % Bottom panel - Metrics
        \draw[thick] (0.5,0.5) rectangle (13.5,4);
        \node at (7, 4.3) {\small Métricas e Controles};
        
        % Buttons
        \draw[fill=blue!20] (1,1) rectangle (3,1.5);
        \node at (2, 1.25) {\tiny Start};
        
        \draw[fill=red!20] (3.5,1) rectangle (5.5,1.5);
        \node at (4.5, 1.25) {\tiny Stop};
        
        \draw[fill=green!20] (6,1) rectangle (8,1.5);
        \node at (7, 1.25) {\tiny Save Results};
        
        % Metrics display
        \node[align=left] at (2.5, 3) {\tiny FPS: 24.3};
        \node[align=left] at (2.5, 2.5) {\tiny MPJPE: 42.1mm};
        \node[align=left] at (2.5, 2) {\tiny Latency: 41ms};
        
    \end{tikzpicture}
    \caption{Layout da interface gráfica de visualização.}
    \label{fig:gui}
\end{figure}

\subsection{Integração entre Módulos}

\subsubsection{Pipeline End-to-End}

\begin{lstlisting}[language=Python, caption=Pipeline Completo]
class FullBodyPoseEstimation:
    def __init__(self, config):
        # Inicializar modulos
        self.camera = CameraAcquisition(
            camera_id=config.camera_id
        )
        self.detector_2d = Keypoint2DDetector().to('cuda')
        self.detector_2d.load_state_dict(
            torch.load(config.detector_weights)
        )
        self.lifting_3d = Lifting3DNetwork().to('cuda')
        self.lifting_3d.load_state_dict(
            torch.load(config.lifting_weights)
        )
        self.visualizer = VisualizationModule()
        
        # Modo de operacao
        self.mode = 'body'  # ou 'face', 'hands', 'full'
        
    def run_inference(self):
        while True:
            # 1. Capturar frame
            batch, temporal_buffer = self.camera.get_batch()
            
            # 2. Detectar keypoints 2D
            start_time = time.time()
            with torch.no_grad():
                outputs_2d = self.detector_2d(
                    batch.cuda(), mode=self.mode
                )
            
            keypoints_2d = outputs_2d['body']  # [B, 17, 3]
            
            # 3. Lifting para 3D
            # Preparar sequencia temporal
            kpts_seq = self._prepare_sequence(
                keypoints_2d, temporal_buffer
            )
            
            with torch.no_grad():
                keypoints_3d = self.lifting_3d(kpts_seq)
            
            inference_time = time.time() - start_time
            
            # 4. Visualizar
            self.visualizer.render_2d(
                batch[0].cpu().numpy(),
                keypoints_2d[0, :, :2].cpu().numpy(),
                keypoints_2d[0, :, 2].cpu().numpy()
            )
            self.visualizer.render_3d(
                keypoints_3d[0].cpu().numpy()
            )
            
            if self.ground_truth_available:
                gt_3d = self.load_ground_truth()
                self.visualizer.update_metrics(
                    keypoints_3d[0].cpu().numpy(),
                    gt_3d,
                    inference_time
                )
            
            plt.pause(0.001)
            
            # Verificar condicao de parada
            if self._should_stop():
                break
    
    def evaluate_on_dataset(self, dataset_path):
        # Avaliacao completa no Drive&Act
        results = {
            'mpjpe': [],
            'pa_mpjpe': [],
            'pck_50': [],
            'fps': []
        }
        
        dataloader = self._create_dataloader(dataset_path)
        
        for batch in tqdm(dataloader):
            frames, gt_2d, gt_3d = batch
            
            start = time.time()
            pred_3d = self._inference_pipeline(frames)
            elapsed = time.time() - start
            
            # Calcular metricas
            mpjpe = compute_mpjpe(pred_3d, gt_3d)
            pa_mpjpe = compute_pa_mpjpe(pred_3d, gt_3d)
            pck = compute_pck(pred_3d, gt_3d, threshold=50)
            fps = len(frames) / elapsed
            
            results['mpjpe'].append(mpjpe)
            results['pa_mpjpe'].append(pa_mpjpe)
            results['pck_50'].append(pck)
            results['fps'].append(fps)
        
        # Agregar resultados
        final_results = {
            'mpjpe_mean': np.mean(results['mpjpe']),
            'mpjpe_std': np.std(results['mpjpe']),
            'pa_mpjpe_mean': np.mean(results['pa_mpjpe']),
            'pck_50_mean': np.mean(results['pck_50']),
            'fps_mean': np.mean(results['fps'])
        }
        
        return final_results
\end{lstlisting}

% --- CRONOGRAMA ---
\section{CRONOGRAMA}

O cronograma de implementação está dividido em 6 fases ao longo de 11 meses, seguindo a estratégia incremental proposta.

\begin{table}[H]
    \centering
    \caption{Cronograma detalhado de implementação.}
    \label{tab:cronograma_impl}
    \small
    \begin{tabular}{|c|p{5.5cm}|c|p{2.8cm}|c|}
        \hline
        \textbf{Fase} & \textbf{Atividades} & \textbf{Período} & \textbf{Entrega} & \textbf{Status} \\ \hline
        
        1 & \textbf{Setup e Corpo 2D}
        \newline - Configuração ambiente
        \newline - Módulo 1: Aquisição
        \newline - Módulo 2: Detector corpo 2D
        \newline - Treinamento em COCO grayscale & 
        set--out/2025 & 
        Pipeline corpo 2D funcional & 
        \textcolor{orange}{Em andamento} \\ \hline
        
        2 & \textbf{Lifting 3D Corporal}
        \newline - Módulo 3: Lifting network
        \newline - Treinamento em Human3.6M
        \newline - Fine-tuning Drive\&Act
        \newline - Módulo 4: Visualização básica &
        nov--dez/2025 &
        Sistema corpo 3D completo &
        Não iniciado \\ \hline
        
        3 & \textbf{Documentação e Paper}
        \newline - Escrita do artigo
        \newline - Experimentos adicionais
        \newline - Comparação com baselines
        \newline - Submissão &
        jan/2026 &
        \textbf{Paper submetido} &
        Não iniciado \\ \hline
        
        4 & \textbf{Integração Face}
        \newline - Implementar face\_head
        \newline - Treinar em 300W grayscale
        \newline - Integrar lifting facial
        \newline - Testes de direção olhar &
        fev--mar/2026 &
        Módulo face integrado &
        Não iniciado \\ \hline
        
        5 & \textbf{Integração Mãos}
        \newline - Implementar hands\_head
        \newline - Treinar em FreiHAND
        \newline - Sistema full-body unificado
        \newline - Otimização fim-a-fim &
        abr--mai/2026 &
        Sistema full-body completo &
        Não iniciado \\ \hline
        
        6 & \textbf{Finalização}
        \newline - Otimização (quantização)
        \newline - Bateria completa de testes
        \newline - Documentação técnica
        \newline - Escrita do TCC
        \newline - Preparação apresentação &
        jun--ago/2026 &
        \textbf{TCC completo} &
        Não iniciado \\ \hline
    \end{tabular}
\end{table}

\subsection{Marcos de Validação}

\begin{itemize}
    \item \textbf{Marco 1 (Out/2025)}: Detector 2D corporal operando em imagens IR do Drive\&Act com PCK@50mm $> 85\%$
    \item \textbf{Marco 2 (Dez/2025)}: Sistema corpo 3D com MPJPE $< 70$mm no Drive\&Act
    \item \textbf{Marco 3 (Jan/2026)}: Paper aceito ou em revisão em conferência relevante
    \item \textbf{Marco 4 (Mar/2026)}: Sistema corpo + face com detecção de olhar funcional
    \item \textbf{Marco 5 (Mai/2026)}: Sistema full-body operando a $\geq 15$ FPS em RTX 3060
    \item \textbf{Marco 6 (Ago/2026)}: TCC completo e aprovado pelos orientadores
\end{itemize}

% --- PROCEDIMENTOS DE TESTE E VALIDAÇÃO ---
\section{PROCEDIMENTOS DE TESTE E VALIDAÇÃO}

\subsection{Testes em Caixa Preta (Perspectiva do Usuário)}

Os testes em caixa preta avaliarão o sistema do ponto de vista do usuário final, sem conhecimento da implementação interna.

\subsubsection{Teste 1: Precisão da Estimação}

\textbf{Objetivo:} Validar se o sistema detecta corretamente a pose do ocupante.

\noindent\textbf{Procedimento:}
\begin{enumerate}
    \item Capturar sequências de vídeo de um ocupante em diferentes posturas (normal, curvado, braços estendidos, etc.)
    \item Executar o sistema para estimar pose 3D
    \item Comparar com ground truth do Drive\&Act
    \item Calcular MPJPE médio
\end{enumerate}

\noindent\textbf{Critério de Aceite:} MPJPE $< 50$mm em 90\% dos casos

\subsubsection{Teste 2: Desempenho em Tempo Real}

\textbf{Objetivo:} Verificar se o sistema opera em tempo real.

\noindent\textbf{Procedimento:}
\begin{enumerate}
    \item Processar 1000 frames consecutivos
    \item Medir tempo de processamento de cada frame
    \item Calcular FPS médio e latência
\end{enumerate}

\noindent\textbf{Critério de Aceite:} FPS $\geq 20$, latência $< 100$ms

\subsubsection{Teste 3: Robustez a Oclusões}

\textbf{Objetivo:} Avaliar comportamento com oclusões típicas.

\noindent\textbf{Procedimento:}
\begin{enumerate}
    \item Selecionar frames do Drive\&Act com oclusões de volante
    \item Avaliar degradação de precisão
    \item Verificar se sistema mantém tracking temporal
\end{enumerate}

\noindent\textbf{Critério de Aceite:} MPJPE $< 75$mm mesmo com oclusões parciais

\subsubsection{Teste 4: Variação de Iluminação}

\textbf{Objetivo:} Verificar robustez a diferentes condições de luz.

\noindent\textbf{Procedimento:}
\begin{enumerate}
    \item Testar em subsets do Drive\&Act: dia, noite, túnel
    \item Comparar métricas entre condições
\end{enumerate}

\noindent\textbf{Critério de Aceite:} Variação de MPJPE $< 20\%$ entre condições

\subsection{Testes em Caixa Branca (Perspectiva do Sistema)}

Os testes em caixa branca avaliarão componentes internos do sistema.

\subsubsection{Teste 5: Módulo de Aquisição}

\textbf{Verificações:}
\begin{itemize}
    \item Correção de distorção funciona corretamente (comparar com calibração OpenCV)
    \item Normalização gera valores em [0, 1]
    \item Buffer temporal mantém 16 frames mais recentes
    \item Sincronização entre múltiplas câmeras (se aplicável)
\end{itemize}

\subsubsection{Teste 6: Detector 2D}

\textbf{Verificações:}
\begin{itemize}
    \item Heatmaps têm picos únicos e bem definidos
    \item Coordenadas extraídas correspondem aos picos
    \item Confiança correlaciona com qualidade da detecção
    \item Teste unitário: imagens sintéticas com keypoints conhecidos
\end{itemize}

\subsubsection{Teste 7: Lifting 3D}

\textbf{Verificações:}
\begin{itemize}
    \item Coordenadas 3D respeitam restrições biomecânicas (comprimentos de membros)
    \item Fusão temporal reduz jitter (comparar com single-frame)
    \item Teste com keypoints 2D ground truth: avaliar capacidade pura de lifting
    \item Análise de atenção: verificar se modelo foca em relações relevantes
\end{itemize}

\subsubsection{Teste 8: Pipeline Integrado}

\textbf{Verificações:}
\begin{itemize}
    \item Propagação correta de dados entre módulos
    \item Gerenciamento de memória GPU (sem vazamentos)
    \item Tratamento de exceções e casos limite
    \item Profiling de performance: identificar gargalos
\end{itemize}

\subsection{Validação no Dataset Drive\&Act}

\textbf{Protocol de Validação Completa:}

\begin{lstlisting}[language=Python, caption=Script de Validação]
def validate_on_driveact(model, split='test'):
    dataset = DriveActDataset(split=split)
    results = defaultdict(list)
    
    for sample in tqdm(dataset):
        frame_ir, gt_kpts_2d, gt_kpts_3d, metadata = sample
        
        # Inferencia
        pred_kpts_2d = model.detect_2d(frame_ir)
        pred_kpts_3d = model.lift_to_3d(pred_kpts_2d)
        
        # Metricas
        mpjpe = compute_mpjpe(pred_kpts_3d, gt_kpts_3d)
        pa_mpjpe = compute_pa_mpjpe(pred_kpts_3d, 
                                     gt_kpts_3d)
        pck_50 = compute_pck(pred_kpts_3d, gt_kpts_3d, 50)
        
        # Agrupar por condicao
        condition = metadata['illumination']  
        # 'day', 'night', 'tunnel'
        results[condition].append({
            'mpjpe': mpjpe,
            'pa_mpjpe': pa_mpjpe,
            'pck_50': pck_50
        })
    
    # Reportar resultados agregados
    report = generate_validation_report(results)
    return report
\end{lstlisting}

\textbf{Métricas Reportadas:}
\begin{itemize}
    \item MPJPE (mm): Erro médio por junta
    \item PA-MPJPE (mm): Erro após alinhamento de Procrustes
    \item PCK@50mm (\%): Porcentagem de keypoints corretos
    \item FPS: Taxa de processamento
    \item Latência (ms): Tempo end-to-end
    \item Uso de memória GPU (MB)
\end{itemize}

% --- ANÁLISE DE RISCOS ---
\section{RISCOS}

\subsection{Riscos Técnicos}

\begin{table}[H]
    \centering
    \caption{Análise de riscos técnicos.}
    \small
    \begin{tabular}{|p{3cm}|c|c|c|p{4cm}|}
        \hline
        \textbf{Risco} & \textbf{Prob.} & \textbf{Impacto} & \textbf{Sever.} & \textbf{Mitigação} \\ \hline
        
        Complexidade integração full-body & 
        Alta & 
        Alto & 
        9 & 
        Desenvolvimento incremental; validação por etapas; baselines separadas \\ \hline
        
        Não atingir tempo real & 
        Média & 
        Alto & 
        6 & 
        Arquiteturas eficientes desde início; quantização; pruning; aceitar FPS reduzido (15 FPS mínimo) \\ \hline
        
        Transferência domínio inadequada & 
        Média & 
        Médio & 
        4 & 
        Fine-tuning progressivo; data augmentation específica; coleta dados complementares \\ \hline
        
        Overfitting em Drive\&Act & 
        Baixa & 
        Médio & 
        2 & 
        Validação cruzada; regularização; early stopping \\ \hline
        
    \end{tabular}
\end{table}

\textbf{Ações de Contingência:}
\begin{itemize}
    \item Se integração full-body for inviável: entregar sistema corpo + face (mãos como futuro trabalho)
    \item Se tempo real não for atingido: focar em precisão e documentar trade-offs; sugerir hardware mais potente
    \item Se Drive\&Act tiver qualidade insuficiente: complementar com dados sintéticos renderizados
\end{itemize}

\subsection{Riscos de Dados}

\begin{table}[H]
    \centering
    \caption{Análise de riscos de dados.}
    \small
    \begin{tabular}{|p{3cm}|c|c|c|p{4cm}|}
        \hline
        \textbf{Risco} & \textbf{Prob.} & \textbf{Impacto} & \textbf{Sever.} & \textbf{Mitigação} \\ \hline
        
        Anotações Drive\&Act com ruído & 
        Média & 
        Médio & 
        4 & 
        Análise qualitativa; filtragem de outliers; loss functions robustas \\ \hline
        
        Cobertura limitada de variações & 
        Média & 
        Médio & 
        4 & 
        Identificação precoce de gaps; dados sintéticos; augmentation agressiva \\ \hline
        
        Acesso limitado a GPUs & 
        Baixa & 
        Alto & 
        3 & 
        Uso de Google Colab; clusters universitários; cloud (AWS, Azure) \\ \hline
        
    \end{tabular}
\end{table}

\subsection{Riscos de Cronograma}

\begin{table}[H]
    \centering
    \caption{Análise de riscos de cronograma.}
    \small
    \begin{tabular}{|p{3cm}|c|c|c|p{4cm}|}
        \hline
        \textbf{Risco} & \textbf{Prob.} & \textbf{Impacto} & \textbf{Sever.} & \textbf{Mitigação} \\ \hline
        
        Atraso na Fase 1 (corpo 2D) & 
        Média & 
        Alto & 
        6 & 
        Buffer de 2 semanas; priorização clara; orientação semanal \\ \hline
        
        Paper não aceito na 1ª submissão & 
        Média & 
        Baixo & 
        2 & 
        Continuar desenvolvimento; resubmeter em paralelo; não bloquear progresso \\ \hline
        
        Problemas pessoais/saúde & 
        Baixa & 
        Alto & 
        3 & 
        Comunicação imediata com orientador; replanejar cronograma \\ \hline
        
    \end{tabular}
\end{table}

\textbf{Monitoramento de Riscos:} Revisão quinzenal em reuniões com orientador; atualização da tabela de riscos; acionamento de contingências quando severidade > 6.

% --- CONCLUSÃO ---
\section{CONCLUSÃO}

Este Projeto Físico apresenta a especificação completa e detalhada do sistema de estimação de pose full-body para ambiente veicular. O projeto está estruturado em quatro módulos principais claramente definidos: Aquisição e Pré-processamento, Detecção 2D de Keypoints, Fusão Temporal e Lifting 3D, e Visualização e Métricas.

As principais tecnologias selecionadas - Python, PyTorch, OpenCV, arquiteturas baseadas em RTMPose e Transformers - são amplamente utilizadas na comunidade acadêmica e industrial, garantindo suporte, documentação e recursos disponíveis. A escolha de PyTorch como framework principal permite flexibilidade no desenvolvimento e facilita a implementação de arquiteturas customizadas.

A estratégia de desenvolvimento incremental, iniciando com pose corporal (Fase 1-2), seguida de publicação científica (Fase 3), e então expandindo para face (Fase 4) e mãos (Fase 5), oferece múltiplas vantagens. Primeiro, permite validação cuidadosa de cada componente antes da integração. Segundo, reduz riscos ao garantir entregas parciais de valor mesmo se o sistema completo não for concluído no prazo. Terceiro, possibilita feedback da comunidade científica através da publicação antes da finalização completa do TCC.

Os procedimentos de teste e validação foram cuidadosamente planejados em duas perspectivas complementares. Os testes em caixa preta avaliarão o sistema do ponto de vista do usuário final, focando em métricas práticas como precisão (MPJPE < 50mm), desempenho em tempo real ($\geq$ 20 FPS), robustez a oclusões e variações de iluminação. Os testes em caixa branca validarão componentes internos, garantindo que cada módulo funciona corretamente de forma isolada e integrada. A validação no dataset Drive\&Act, específico para contexto automotivo, fornecerá métricas realistas de performance em cenários de aplicação real.

A análise de riscos identificou três categorias principais: técnicos, de dados e de cronograma. Os riscos técnicos mais significativos incluem a complexidade da integração full-body (severidade 9) e a possibilidade de não atingir processamento em tempo real (severidade 6). Para cada risco identificado, estratégias de mitigação específicas foram definidas, incluindo desenvolvimento incremental, uso de arquiteturas eficientes, e ações de contingência claramente estabelecidas. Por exemplo, se a integração full-body se mostrar inviável no prazo, o sistema será entregue com corpo mais face, mantendo mãos como trabalho futuro.

Atualmente, o projeto é considerado viável com base em múltiplas evidências. Primeiro, trabalhos relacionados como RTMPose e LiftPose3D demonstram que as técnicas propostas são factíveis. Segundo, o dataset Drive\&Act fornece dados específicos do domínio para treinamento e validação. Terceiro, recursos computacionais adequados (GPUs NVIDIA de gama média) estão disponíveis através de clusters universitários e plataformas cloud. Quarto, a estratégia incremental e as contingências planejadas garantem flexibilidade para adaptar o escopo se necessário.

Os riscos atualmente avaliados como mais críticos (severidade $\geq$ 6) serão monitorados quinzenalmente em reuniões com o orientador. A complexidade da integração full-body será endereçada através de validação rigorosa de cada modalidade separadamente antes da unificação. O risco de não atingir tempo real será mitigado através de otimizações progressivas (quantização, pruning) e ajuste de expectativas (aceitar 15 FPS como mínimo aceitável para prova de conceito).

O cronograma estabelecido, distribuído em 11 meses com 6 fases claramente definidas, é realista e foi elaborado considerando a experiência prévia do orientador e trabalhos similares na literatura. Os marcos de validação intermediários (6 marcos ao longo do projeto) permitirão identificação precoce de desvios e acionamento de ações corretivas.

Em relação à evolução dos riscos desde o planejamento inicial, observa-se que a disponibilidade do dataset Drive\&Act reduziu significativamente a incerteza sobre dados de treinamento, diminuindo a probabilidade de riscos relacionados. Por outro lado, a compreensão mais profunda da complexidade da fusão temporal aumentou ligeiramente a avaliação de impacto deste componente, reforçando a necessidade de atenção especial durante sua implementação.

A equipe (discente e orientadores) avalia que o projeto continua sendo plenamente viável e com grande potencial de contribuição científica e prática. A combinação de um problema relevante, solução tecnicamente sólida, recursos adequados e planejamento cuidadoso fornece confiança de que os objetivos serão alcançados. O sistema proposto, uma vez concluído, poderá contribuir tanto para o avanço do estado da arte em estimação de pose em ambientes restritos quanto para aplicações práticas em segurança veicular, potencialmente beneficiando milhões de usuários através da democratização de tecnologias de monitoramento de ocupantes.

A documentação técnica detalhada apresentada neste Projeto Físico - incluindo pseudocódigo, diagramas, especificações de interfaces e estratégias de treinamento - demonstra que não existem lacunas significativas no entendimento de como o sistema será implementado. Todos os componentes críticos foram especificados em nível de detalhe suficiente para iniciar a implementação, e as ferramentas, bibliotecas e frameworks necessários foram identificados e estão disponíveis.

% --- REFERÊNCIAS BIBLIOGRÁFICAS ---
\section{REFERÊNCIAS BIBLIOGRÁFICAS}

\begin{thebibliography}{99}

\bibitem{pytorch}
PASZKE, A. et al. \textbf{PyTorch: An Imperative Style, High-Performance Deep Learning Library}. In: Advances in Neural Information Processing Systems (NeurIPS), 2019. p. 8024-8035. Disponível em: \url{https://pytorch.org}. Acesso em: 22 out. 2025.

\bibitem{opencv}
BRADSKI, G. \textbf{The OpenCV Library}. Dr. Dobb's Journal of Software Tools, v. 25, n. 11, p. 120-125, 2000. Disponível em: \url{https://opencv.org}. Acesso em: 22 out. 2025.

\bibitem{rtmpose}
JIANG, T. et al. \textbf{RTMPose: Real-Time Multi-Person Pose Estimation based on MMPose}. arXiv preprint arXiv:2303.07399, 2023. Disponível em: \url{https://arxiv.org/abs/2303.07399}. Acesso em: 22 out. 2025.

\bibitem{martinez2017}
MARTINEZ, J. et al. \textbf{A simple yet effective baseline for 3D human pose estimation}. In: IEEE International Conference on Computer Vision (ICCV), 2017. p. 2640-2649. DOI: 10.1109/ICCV.2017.288.

\bibitem{iskakov2019}
ISKAKOV, K. et al. \textbf{Learnable triangulation of human pose}. In: IEEE International Conference on Computer Vision (ICCV), 2019. p. 7718-7727. DOI: 10.1109/ICCV.2019.00781.

\bibitem{driveact}
MARTIN, M. et al. \textbf{Drive\&Act: A Multi-Modal Dataset for Fine-Grained Driver Behavior Recognition in Autonomous Vehicles}. In: IEEE International Conference on Computer Vision (ICCV), 2019. p. 2801-2810. DOI: 10.1109/ICCV.2019.00289.

\bibitem{vaswani2017}
VASWANI, A. et al. \textbf{Attention is all you need}. In: Advances in Neural Information Processing Systems (NeurIPS), 2017. p. 5998-6008. Disponível em: \url{https://arxiv.org/abs/1706.03762}. Acesso em: 22 out. 2025.

\bibitem{liftpose3d}
ZHANG, C. et al. \textbf{LiftPose3D: Transforming 2D poses to 3D in real-time}. arXiv preprint arXiv:2008.02945, 2020. Disponível em: \url{https://arxiv.org/abs/2008.02945}. Acesso em: 22 out. 2025.

\bibitem{coco}
LIN, T.-Y. et al. \textbf{Microsoft COCO: Common Objects in Context}. In: European Conference on Computer Vision (ECCV), 2014. p. 740-755. DOI: 10.1007/978-3-319-10602-1\_48.

\bibitem{mpii}
ANDRILUKA, M. et al. \textbf{2D Human Pose Estimation: New Benchmark and State of the Art Analysis}. In: IEEE Conference on Computer Vision and Pattern Recognition (CVPR), 2014. p. 3686-3693. DOI: 10.1109/CVPR.2014.471.

\bibitem{human36m}
IONESCU, C. et al. \textbf{Human3.6M: Large Scale Datasets and Predictive Methods for 3D Human Sensing in Natural Environments}. IEEE Transactions on Pattern Analysis and Machine Intelligence, v. 36, n. 7, p. 1325-1339, 2014. DOI: 10.1109/TPAMI.2013.248.

\bibitem{efficientnet}
TAN, M.; LE, Q. V. \textbf{EfficientNet: Rethinking Model Scaling for Convolutional Neural Networks}. In: International Conference on Machine Learning (ICML), 2019. p. 6105-6114. Disponível em: \url{https://arxiv.org/abs/1905.11946}. Acesso em: 22 out. 2025.

\bibitem{timm}
WIGHTMAN, R. \textbf{PyTorch Image Models (timm)}. GitHub repository, 2019. Disponível em: \url{https://github.com/rwightman/pytorch-image-models}. Acesso em: 22 out. 2025.

\bibitem{mmpose}
CONTRIBUTORS, M. \textbf{OpenMMLab Pose Estimation Toolbox and Benchmark}. GitHub repository, 2020. Disponível em: \url{https://github.com/open-mmlab/mmpose}. Acesso em: 22 out. 2025.

\bibitem{300w}
SAGONAS, C. et al. \textbf{300 Faces in-the-Wild Challenge: The First Facial Landmark Localization Challenge}. In: IEEE International Conference on Computer Vision Workshops (ICCVW), 2013. p. 397-403. DOI: 10.1109/ICCVW.2013.59.

\bibitem{freihand}
ZIMMERMANN, C. et al. \textbf{FreiHAND: A Dataset for Markerless Capture of Hand Pose and Shape from Single RGB Images}. In: IEEE International Conference on Computer Vision (ICCV), 2019. p. 813-822. DOI: 10.1109/ICCV.2019.00090.

\bibitem{adam}
KINGMA, D. P.; BA, J. \textbf{Adam: A Method for Stochastic Optimization}. In: International Conference on Learning Representations (ICLR), 2015. Disponível em: \url{https://arxiv.org/abs/1412.6980}. Acesso em: 22 out. 2025.

\bibitem{batchnorm}
IOFFE, S.; SZEGEDY, C. \textbf{Batch Normalization: Accelerating Deep Network Training by Reducing Internal Covariate Shift}. In: International Conference on Machine Learning (ICML), 2015. p. 448-456. Disponível em: \url{https://arxiv.org/abs/1502.03167}. Acesso em: 22 out. 2025.

\bibitem{wandb}
BIEWALD, L. \textbf{Experiment Tracking with Weights and Biases}. Software available from wandb.com, 2020. Disponível em: \url{https://www.wandb.com}. Acesso em: 22 out. 2025.

\bibitem{lightning}
FALCON, W. et al. \textbf{PyTorch Lightning}. GitHub repository, 2019. Disponível em: \url{https://github.com/Lightning-AI/lightning}. Acesso em: 22 out. 2025.

\end{thebibliography}

\end{document}